\documentclass[a4paper,12pt, headsepline]{scrartcl}

\usepackage[utf8]{inputenc}
\usepackage[T1]{fontenc}
\pagestyle{plain}
\usepackage{lmodern}
\usepackage[english]{babel}
\usepackage[official]{eurosym}
\usepackage{ntheorem}
\newtheorem{hyp}{Hypothesis}
\usepackage{pdfpages}
\usepackage{adjustbox}
\usepackage{footnote}
\usepackage{titling}
\usepackage{csquotes}
\usepackage[objectset = centering]{floatrow}
\usepackage{moreverb}
\usepackage{varioref}
\usepackage{setspace}
\usepackage{graphicx}
\usepackage{tabularx}
\usepackage{dcolumn}
\usepackage{color}
\usepackage{amsmath} 
\usepackage{amsfonts}
\usepackage{mathtools}
\usepackage{relsize}
\usepackage{bm}
\usepackage{dsfont}
\usepackage{algorithm2e}
\usepackage{enumerate}
\usepackage{float}
\usepackage[mathscr]{euscript}
\usepackage[
colorlinks = true,
linkcolor = black,
citecolor = blue,
urlcolor = teal 
]{hyperref}

\setlength{\parindent}{0ex} 
\setlength{\parskip}{1ex}


\usepackage[automark]{scrlayer-scrpage}
\setkomafont{pagehead}{\scshape}
\pagestyle{scrheadings}
\ihead{\headmark}
\chead{}
\ohead{}

\cfoot{\pagemark}

\usepackage{csquotes}
\usepackage[
backend=biber,
natbib=true,
language = english,
doi = false, url = false, isbn = false, eprint = false,
style = apa]
{biblatex}
%%
\usepackage{booktabs}
\usepackage{longtable}
\usepackage{array}
\usepackage{multirow}
\usepackage{wrapfig}
\usepackage{colortbl}
\usepackage{pdflscape}
% \usepackage{tabu} % removed: tabu is unmaintained and no longer shipped by MiKTeX; use booktabs/tabularx/longtable instead
% \usepackage{threeparttable}
% \usepackage{threeparttablex}
% \usepackage[normalem]{ulem}
\usepackage{makecell}
\usepackage{xcolor}
%%

\DeclareLanguageMapping{english}{english-apa}
\addbibresource{literature.bib}

%% Commands %%
\DeclareMathOperator*{\argminA}{arg\,min}
\DeclareMathOperator*{\argmaxA}{arg\,max}
\DeclareMathOperator*{\maxA}{max}
\DeclareMathOperator*{\minA}{min}
\newenvironment{proof}{\paragraph{Proof:}}{\hfill$\square$}
%%%%%%%%%%%%%%

% Schrift
\setkomafont{sectioning}{\rmfamily\bfseries\boldmath}

\setcapindent{0em} % kein Einrücken der Caption von Figures und Tabellen
\setcapwidth{0.9\textwidth} % Breite der Caption nur 90% der Textbreite, damit sie sich vom restlichen Text abhebt
\setlength{\abovecaptionskip}{0.2cm} % Abstand der zwischen Bild- und Bildunterschrift
%\renewcommand{\baselinestretch}{1.5}
\numberwithin{equation}{section}
\title{Monetary Policy Transmission and Building Land Scarcity:\\ A Panel Local Projections Analysis of Regional Austrian Housing Markets}
\date{\today}
\author{Lea Schwabe\\(11908829)\\[0.47cm]Supervisor:\\Assoz.~Prof.~Mag.~Dr.~Paul Pichler}
\begin{document}
\begin{titlingpage}
\maketitle
\begin{abstract}
\begin{center}
	\textbf{Abstract}
\end{center}
\vspace*{2pt}
% Abstract text goes here.
\end{abstract}
\end{titlingpage}
\newpage
\thispagestyle{empty}
\vspace*{6.5cm}
\begin{center}
\textbf{Acknowledgements}
\end{center}
\vspace*{-17pt}
% Acknowledgements text goes here.
\newpage
\tableofcontents
\thispagestyle{empty}
\clearpage
\pagenumbering{arabic} 

% =============================================================================
%  1. INTRODUCTION
%  Purpose: Motivate the thesis, define the empirical problem, state the
%  (expected) contribution.
%  Guiding questions: Why is the reaction of regional real estate markets to
%  aggregate monetary policy shocks heterogeneous? Why is the district-level
%  analysis analytically needed? What is the direct contribution versus
%  existing literature in Austria?
% =============================================================================
\section{Introduction}\label{sec:int}

\subsection{Motivation}\label{subsec:motivation}
% Austrian real estate price boom, affordability, monetary policy phases,
% regional differences -- why are regional responses to common shocks
% heterogeneous?

\subsection{Research Gap}\label{subsec:gap}
% National trends are documented; the interaction between ECB shocks and land
% scarcity has not been directly estimated.

\subsection{Research Question}\label{subsec:rq}
% Primary RQ: To what extent do local supply constraints moderate the dynamic
% response of district-level real estate prices to ECB monetary shocks in
% Austria?
% Sub-questions: (a) magnitude of cumulative log price changes across
% elastic/inelastic districts; (b) temporal response over a three-year horizon;
% (c) asymmetry across monetary policy phases.

\subsection{Contribution}\label{subsec:contribution}
% Evidence at the district level.

\subsection{Thesis Outline}\label{subsec:outline}
% Brief roadmap of the remaining chapters.

% =============================================================================
%  2. INSTITUTIONAL AND ECONOMIC BACKGROUND
%  Purpose: Explain the Austrian and Euro Area context (pre literature review).
%  Guiding questions: Which phases of monetary policy (during the sample period)
%  are relevant for Austria? How do districts differ in prices, urbanity and
%  building land? How do districts differ structurally from each other?
% =============================================================================
\section{Institutional and Economic Background}\label{sec:background}

\subsection{ECB Monetary Policy Context}\label{subsec:ecb-context}
% APP, PEPP, rate-hiking cycle; why high-frequency monetary shocks are
% preferable to policy-period dummies.

\subsection{The Austrian Real Estate Market, 2017--2025}\label{subsec:at-market}
% National price trends, descriptive narrative of real estate market
% developments; possibly include mortgage rates during monetary policy phases.

\subsection{Regional Housing Market Heterogeneity}\label{subsec:regional-heterogeneity}
% Why and how do Austrian districts differ structurally: urban/rural
% differences, land scarcity, topography, zoning, population/demographic
% pressures.

% =============================================================================
%  3. LITERATURE REVIEW: THEORY AND HYPOTHESES
%  Purpose: Position the thesis in the monetary transmission literature.
%  Guiding questions: Which mechanisms indicate that inelastic districts react
%  differently to monetary policy shocks? What is the precise research gap for
%  Austria? Which testable hypotheses can be derived from existing literature?
% =============================================================================
\section{Literature Review: Theory and Hypotheses}\label{sec:literature}

\subsection{Monetary Policy and Housing Prices}\label{subsec:mp-housing}
% Housing channel of monetary policy; monetary policy transmission literature.

\subsection{Housing Supply Constraints and Monetary Policy Transmission}\label{subsec:supply-transmission}
% How do regional supply constraints shape monetary policy transmission through
% housing markets? (Possibly a sub-chapter of 3.1.)

\subsection{Austrian Housing and Land Price Evidence}\label{subsec:at-evidence}
% Descriptive literature for Austria (from the proposal).

\subsection{Remaining Gap and Hypotheses}\label{subsec:hypotheses}
% Existing work supports the mechanism, but not yet with Austrian district-level
% price data using local land-scarcity classifications. Derive testable
% hypotheses here.
%
% \begin{hyp}
% 	Stronger price responses in supply-inelastic districts.
% \end{hyp}
% \begin{hyp}
% 	More persistent effects over time in supply-inelastic districts.
% \end{hyp}
% \begin{hyp}
% 	Asymmetric responses across monetary policy phases (smaller reversal
% 	effects in supply-inelastic districts).
% \end{hyp}

% =============================================================================
%  4. DATA AND VARIABLES
%  Purpose: Define the spatial and temporal unit (district level, number of
%  districts, sample period, balanced/unbalanced panel, etc.).
%  Guiding questions: Which variables enter the panel? How is inelasticity
%  measured? How are daily ECB shocks transformed and merged with the quarterly
%  panel? How are districts temporally harmonized?
% =============================================================================
\section{Data and Variables}\label{sec:data}

\subsection{Real Estate Price Data}\label{subsec:price-data}
% Units, data description, price measure, missingness, transformations.

\subsection{Monetary Policy Shocks}\label{subsec:shock-data}
% EA-MPD source, chosen shock variable (OIS 20Y), event-to-quarter aggregation,
% sign convention, scaling.

\subsection{Supply-Constraint Measure}\label{subsec:supply-measure}
% OeROK building land reserves; construction of the inelastic dummy.

\subsection{Control Variables}\label{subsec:controls}
% Unemployment, population, net income per capita: source, frequency,
% transformations, interpolation where applicable.

\subsection{Final Panel Construction}\label{subsec:panel-construction}
% Merging logic, district IDs, missing values, final observation count.

% =============================================================================
%  5. EMPIRICAL STRATEGY
%  Purpose: Explain the local projections approach and empirical model,
%  identification, and interpretation.
%  Guiding questions: How do panel local projections work? Why and how does beta
%  identify the differential response of inelastic districts? Which assumptions
%  are critical?
% =============================================================================
\section{Empirical Strategy}\label{sec:strategy}

\subsection{Panel Local Projections}\label{subsec:lp}
% Methodological background: how are impulse response functions estimated?

\subsection{Baseline Model and Identification}\label{subsec:baseline}
% Discuss inputs, controls, lagged or not, etc.
%
% The following model is specified for each horizon $h = 1, \ldots, 12$:
% \begin{equation}
% 	y_{i,t+h} - y_{i,t-1} = \alpha_{i,h} + \tau_{t,h}
% 		+ \beta_h \left( \text{Shock}_t \times \text{Inelastic}_i \right)
% 		+ \theta_h X_{i,t} + \varepsilon_{i,t+h}. \label{eq:lp}
% \end{equation}

\subsection{Interpretation of the Coefficient}\label{subsec:coef-interpretation}
% Differential response interpretation of $\beta_h$.

\subsection{Outcome Transformation}\label{subsec:outcome}
% Cumulative log change from $t-1$ to $t+h$ (cf. Jorda, Federal Reserve paper).

\subsection{Fixed Effects}\label{subsec:fe}
% District and time fixed effects; stress that including time FE means the model
% estimates a differential response.

\subsection{Inference and Identification Assumptions}\label{subsec:inference}
% Clustering of standard errors (Driscoll--Kraay or two-way clustering?);
% exogeneity of MP shocks; no omitted district-specific shock correlated with
% inelasticity and price responses; validity of the land-reserve proxy.

% =============================================================================
%  6. DESCRIPTIVE EVIDENCE
%  Purpose: Pre-estimation evidence. Should not argue causality; show that the
%  data contains meaningful regional heterogeneity and motivate the LP analysis.
%  Guiding questions: Are there clear regional patterns before estimation? How
%  large and how persistent are the price reactions? Are differences
%  economically relevant or only statistically significant?
% =============================================================================
\section{Descriptive Evidence}\label{sec:descriptive}
% Summary statistics by property type; district classification summary (number
% of elastic/inelastic districts, including/excluding Vienna); map of building
% land reserves (spatial patterns of supply constraints); national (median)
% price by property type (2017--2025); elastic vs. inelastic prices;
% distribution of cumulative price changes; correlation plot (reserve measure
% vs. price level/growth).

% =============================================================================
%  7. EMPIRICAL RESULTS
% =============================================================================
\section{Empirical Results}\label{sec:results}

\subsection{Baseline Results}\label{subsec:baseline-results}
% House prices as primary outcome; IRF plot.

\subsection{Alternative Responses}\label{subsec:alt-responses}
% Apartments, building land.

\subsection{Robustness Checks}\label{subsec:robustness}
% Vienna treatment, alternative thresholds, continuous reserve measure,
% alternative shock aggregation, shorter horizon, etc.

% =============================================================================
%  8. DISCUSSION
%  Purpose: Interpret results, connect to literature, discuss limitations.
% =============================================================================
\section{Discussion}\label{sec:discussion}

\subsection{Interpretation of Findings}\label{subsec:interpretation}
% Are inelastic districts more sensitive? Is the effect persistent? Does the
% response differ by property type? Relate back to the sub-questions (temporal
% response, etc.).

\subsection{Relation to Literature}\label{subsec:relation-literature}

\subsection{Economic Relevance}\label{subsec:economic-relevance}

\subsection{Policy Implications}\label{subsec:policy}

\subsection{Limitations}\label{subsec:limitations}
% Short sample / availability of price data, supply proxy measurement, Vienna,
% causal interpretation.

% =============================================================================
%  9. CONCLUSION
% =============================================================================
\section{Conclusion}\label{sec:con}
% Restate the research question, summarize data and method, summarize the main
% empirical findings, state the contribution, mention limitations and future
% research possibilities.

\newpage
\printbibliography[heading=bibintoc]

\newpage
\appendix
\section{Appendix}\label{sec:appendix}
% Supplementary tables and figures.

% --- Formatting reference snippets (uncomment and adapt as needed) -----------
% Figure (place image files in the charts/ folder):
% \begin{figure}[H]
% 	\includegraphics[width = 14cm]{charts/example.pdf}
% 	\caption{Example Caption}\label{fig:example}
% \end{figure}
%
% Table with booktabs and shaded cells:
% \begin{table}[H]
% 	\centering
% 	\begin{tabular}[t]{lcc}
% 		\toprule
% 		& Column A & Column B \\
% 		\midrule
% 		\cellcolor{gray!6}{Row 1} & \cellcolor{gray!6}{1} & \cellcolor{gray!6}{2}\\
% 		Row 2 & 3 & 4\\
% 		\bottomrule
% 	\end{tabular}
% 	\caption{Example Table}\label{tab:example}
% \end{table}
% ----------------------------------------------------------------------------

\end{document}
